% GENERATED FILE. Edit canonical lessons, then run npm run generate:books.
% canonical-source-hash: fnv1a64:98324256ca399b3a
% canonical-lessons: ZH-C279-dier-tian, ZH-W279-di, ZH-C279-de-shihou, ZH-C279-turan, ZH-W279-tu, ZH-C279-yixiazi, ZH-C279-reading, ZH-C279-practice

\chapter{The Next Day}
\label{ch:zh-the-next-day}
\hlchaptermodality{279}

\begin{chapteropening}
\textbf{By the end of this chapter:} \emph{I can say what was going on when something happened, and follow a short story.}

Saying what was going on when something happened with …\zh{的时候} and \zh{正在}, suddenly, all at once and the next day, and reading a story of thirteen lines, built on two new written characters: \zh{第} and \zh{突}.
\end{chapteropening}

\section[\texorpdfstring{\zh{第二天} (dì'èr tiān)}{第二天 (dì'èr tiān)}]{\zh{第二天} (dì'èr tiān) — the next day}

\label{lesson:ZH-C279-dier-tian}



\begin{quote}
A new chapter, and it opens on the morning after.

\begin{itemize}
\raggedright
  \item \emph{Recall:} say \zh{先洗手} \zh{然后吃早饭} — \emph{xiān xǐ shǒu, ránhòu chī zǎofàn} — \textbf{R1}, the chapter you have finished
  \item \emph{Recall:} say \emph{jiēzhe}, straight after — \textbf{R2}
  \item \emph{Recall:} say \emph{dìyī}, first — \textbf{R3}, durable
  \item \emph{Recall:} say \zh{明天} — \emph{míngtiān}, tomorrow — \textbf{R3}
\end{itemize}
\end{quote}

\subsection*{You'll want to know: \zh{第二天}}
\textbf{\zh{第二天}} — \emph{dì'èr tiān} — \textquotedblleft{}the next day; the following day\textquotedblright{}.

\begin{quote}
\textbf{\zh{第二天我们去了山上}} \emph{dì'èr tiān wǒmen qù le shān shàng} — The next day we went up the mountain.
\end{quote}

In a story, \zh{第二天} is the day after the one you are telling; \zh{明天} counts from today.

\zh{山上} is \zh{山} with \zh{上}, the short form of \emph{shàngbian}: up on the mountain.

The next lesson puts the first character in your hand. It is the dì of \emph{dìyī}, first.

\begin{cousinweb}
\textbf{\zh{第二天}} (\emph{dì'èr tiān}) is built from \zh{第}, which puts a number in order, \zh{二} 'two' and \zh{天} 'day'. All three are inherited from Middle Chinese, \emph{dejH}, \emph{nyijH} and \emph{then} (Guangyun, in Baxter's transcription); the dictionary glosses the first 'order, sequence'. 'The next day' for the whole is the phrase's own sense: the second day of a story.
\end{cousinweb}

\begin{grammarlens}[title={What you've built}]
\zh{第} before a number makes it an ordinal: \emph{dìyī}, first, is one you say already, and \zh{第二} is second.

\zh{第二天} counts from the day the story is on, so it can follow a day in the past: \textbf{\zh{第二天天气很好}} \emph{dì'èr tiān tiānqì hěn hǎo}.
\end{grammarlens}

\subsection*{Your turn}
\begin{itemize}
\raggedright
  \item \emph{Say it:} \emph{dì'èr tiān}
  \item \emph{Say it:} \emph{dì'èr tiān}, once more
  \item \emph{Say it:} \zh{第二天我们去了山上} — \emph{dì'èr tiān wǒmen qù le shān shàng} — The next day we went up the mountain.
  \item \emph{Say it:} \zh{第二天天气很好} — \emph{dì'èr tiān tiānqì hěn hǎo} — The next day the weather was fine.
  \item \emph{Recall:} say \emph{dìyī}, then \zh{第二天}
\end{itemize}

\begin{tcolorbox}[breakable,colback=teal!4,colframe=teal!35!black,title={Before you move on}]
What does \zh{第二天} mean? (\textquotedblleft{}The next day\textquotedblright{}.) How is it different from \zh{明天}? (\textbf{\zh{明天} counts from today; \zh{第二天} from the day in the story}.)
\end{tcolorbox}
\section[\texorpdfstring{\zh{第} (dì)}{第 (dì)}]{\zh{第} — the bamboo piece of \zh{笑}, then a five-stroke body}

\label{lesson:ZH-W279-di}



\begin{quote}
The word first, then the character inside it.

\begin{itemize}
\raggedright
  \item \emph{Recall:} say \emph{dì'èr tiān} — \textbf{R1}, one lesson back, before it decays
\end{itemize}

\begin{quote}
\zh{笑} \zh{号}
\end{quote}

\emph{Read it:} each one aloud, with nothing said first — \emph{xiào}, \emph{hào}
\end{quote}

\subsection*{Script — the bamboo piece of \zh{笑}, then a five-stroke body}
\hlblockfigure{\detokenize{figures/ZH-W279-di-filmstrip.pdf}}{How \zh{第} is written, stroke by stroke}

\begin{quote}
\zh{第}
\end{quote}

\textbf{Two pieces, and you write the first already.}

\textbf{The bamboo piece, six strokes}, as on top of \zh{笑}: on the left a left-falling stroke, a short horizontal from it and a dot; then the same three on the right.

\textbf{The body, five strokes}: a horizontal that turns down at its right end; a horizontal under it; a stroke that goes down, right, down again and hooks to the left, folding the way the foot of \zh{号} does; the vertical down the middle; and last a left-falling stroke to the lower left.

6 + 5 = eleven. The left-falling stroke comes after the vertical, last of all.

\subsection*{Your turn}
\begin{itemize}
\raggedright
  \item \emph{Write it:} \zh{第}, all eleven, with the model in front of you
  \item \emph{Cover:} write it again from memory, and count the pieces, not the strokes
\end{itemize}

\begin{tcolorbox}[breakable,colback=teal!4,colframe=teal!35!black,title={Before you move on}]
How many strokes in \zh{第}? (\textbf{Eleven}.) Where have you written its top? (\textbf{On \zh{笑}}.) Which stroke comes last? (\textbf{The left-falling stroke, after the vertical}.)

Source: \href{https://raw.githubusercontent.com/chanind/hanzi-writer-data/68d10a4b21150cae5e1ebbd223eed289cf32d90c/data/\%E7\%AC\%AC.json}{Hanzi Writer Data \zh{第}, pinned stroke paths and medians}.
\end{tcolorbox}
\section[\texorpdfstring{…\zh{的时候} (… de shíhou)}{…的时候 (… de shíhou)}]{…\zh{的时候} — when, and what was going on}

\label{lesson:ZH-C279-de-shihou}



\begin{quote}
A word for when, and a word for in the middle of.

\begin{itemize}
\raggedright
  \item \emph{Read it:} \zh{第二天}, aloud — \emph{dì'èr tiān}
  \item \emph{Recall:} say \zh{时候} — \emph{shíhou}, time; when — \textbf{R3}, durable
  \item \emph{Recall:} say \zh{我正在吃饭} — \emph{wǒ zhèngzài chī fàn} — \textbf{R3}
\end{itemize}
\end{quote}

\begin{grammarlens}[title={Grammar Lens: …\zh{的时候} and \zh{正在}, for what was going on}]
\begin{quote}
\textbf{…\zh{的时候} after a clause says when; \zh{正在} in the other clause says what was already going on at that moment.}
\end{quote}

\begin{quote}
\textbf{\zh{他来的时候} \zh{我正在吃饭}} \emph{tā lái de shíhou, wǒ zhèngzài chī fàn} — When he came, I was eating. \textbf{\zh{我回家的时候} \zh{妈妈正在做饭}} \emph{wǒ huí jiā de shíhou, māma zhèngzài zuòfàn} — When I got home, Mum was cooking.
\end{quote}

The clause before \zh{的时候} names the moment, and a bare verb is enough: \zh{他来的时候}. The other clause is the scene that was already running, with the \zh{正在} of \zh{我正在吃饭}.

\zh{正在} has no tense of its own: the story around it puts the scene in the past.

All three characters of \zh{的时候} are ones you write, so it prints whole for the first time.
\end{grammarlens}

\subsection*{Your turn}
Say these aloud:

\begin{itemize}
\raggedright
  \item \zh{他来的时候} \zh{我正在吃饭} — \emph{tā lái de shíhou, wǒ zhèngzài chī fàn} — When he came, I was eating.
  \item \zh{我回家的时候} \zh{妈妈正在做饭} — \emph{wǒ huí jiā de shíhou, māma zhèngzài zuòfàn} — When I got home, Mum was cooking.
  \item \zh{我看书的时候} \zh{他来了} — \emph{wǒ kàn shū de shíhou, tā lái le} — He came while I was reading.
  \item what you were doing when somebody phoned you last
\end{itemize}

\begin{tcolorbox}[breakable,colback=teal!4,colframe=teal!35!black,title={Before you move on}]
What does …\zh{的时候} give you? (\textbf{The moment: when}.) Which word gives the scene that was already going on? (\textbf{\zh{正在}}.)
\end{tcolorbox}
\section[\texorpdfstring{\zh{突然} (tūrán)}{突然 (tūrán)}]{\zh{突然} (tūrán) — suddenly}

\label{lesson:ZH-C279-turan}



\begin{quote}
Before the new one: a moment, and then something breaking into it.

\begin{itemize}
\raggedright
  \item \emph{Read it:} \zh{他来的时候} \zh{我正在吃饭}, aloud — \emph{tā lái de shíhou, wǒ zhèngzài chī fàn}
  \item \emph{Read it:} \zh{然后}, aloud — \emph{ránhòu}
  \item \emph{Recall:} say \zh{下雨} — \emph{xiàyǔ}, to rain — \textbf{R3}, durable
\end{itemize}
\end{quote}

\subsection*{You'll want to know: \zh{突然}}
\textbf{\zh{突然}} — \emph{tūrán} — \textquotedblleft{}suddenly; all of a sudden\textquotedblright{}.

\begin{quote}
\textbf{\zh{突然下雨了}} \emph{tūrán xiàyǔ le} — Suddenly it started to rain.
\end{quote}

\begin{quote}
\textbf{\zh{他突然来了}} \emph{tā tūrán lái le} — He turned up out of the blue.
\end{quote}

The next lesson puts the first character in your hand. The second is the \zh{然} of \zh{然后}.

\begin{cousinweb}
\textbf{\zh{突然}} (\emph{tūrán}) is built from \zh{突} and \zh{然}, both inherited from Middle Chinese. The Guangyun lists \zh{突} under two readings: \emph{thwot}, glossed 'the look of coming out', and \emph{dwot}, glossed 'to butt' and 'to deceive'; tū is the regular modern sound of \emph{thwot}. \zh{然} is \emph{nyen}, glossed as a particle and 'so, like this' (Guangyun, in Baxter's transcription). 'Suddenly' for the pair is the word's own sense, which neither gloss gives.
\end{cousinweb}

\begin{grammarlens}[title={What you've built}]
\zh{突然} goes before the verb, or at the front of the clause: \zh{突然下雨了}, \zh{他突然来了}.

With …\zh{的时候} it breaks into a scene: \textbf{\zh{我看书的时候} \zh{他突然来了}} \emph{wǒ kàn shū de shíhou, tā tūrán lái le}.
\end{grammarlens}

\subsection*{Your turn}
\begin{itemize}
\raggedright
  \item \emph{Say it:} \emph{tūrán}
  \item \emph{Say it:} \emph{tūrán}, once more
  \item \emph{Say it:} \zh{突然下雨了} — \emph{tūrán xiàyǔ le} — Suddenly it started to rain.
  \item \emph{Say it:} \zh{我看书的时候} \zh{他突然来了} — \emph{wǒ kàn shū de shíhou, tā tūrán lái le} — I was reading when he suddenly came.
  \item \emph{Recall:} say \zh{然后}, then \zh{突然}: the same \zh{然}, two words
\end{itemize}

\begin{tcolorbox}[breakable,colback=teal!4,colframe=teal!35!black,title={Before you move on}]
What does \zh{突然} mean? (\textquotedblleft{}Suddenly\textquotedblright{}.) Where does it go? (\textbf{Before the verb, or at the front of its clause}.)
\end{tcolorbox}
\section[\texorpdfstring{\zh{突} (tū)}{突 (tū)}]{\zh{突} — the roof of \zh{家}, \zh{八}, then \zh{大} with a dot}

\label{lesson:ZH-W279-tu}



\begin{quote}
The word first, then the character inside it.

\begin{itemize}
\raggedright
  \item \emph{Recall:} say \emph{tūrán} — \textbf{R1}, one lesson back, before it decays
\end{itemize}

\begin{quote}
\zh{家} \zh{八} \zh{然}
\end{quote}

\emph{Read it:} each one aloud, with nothing said first — \emph{jiā}, \emph{bā}, \emph{rán}
\end{quote}

\subsection*{Script — the roof of \zh{家}, \zh{八}, then \zh{大} with a dot}
\hlblockfigure{\detokenize{figures/ZH-W279-tu-filmstrip.pdf}}{How \zh{突} is written, stroke by stroke}

\begin{quote}
\zh{突}
\end{quote}

\textbf{Three pieces, and you write every one of them already.}

\textbf{The roof, three strokes}, as on top of \zh{家}: the dot, a short left stroke, then a horizontal that hooks down at its right end.

\textbf{\zh{八}, two strokes}, under the roof: a left-falling stroke, then a dot falling to the right.

\textbf{\zh{大} with a dot, four strokes}, at the foot, as on the right of \zh{然}: the bar, the left-falling stroke, the right-falling stroke, and the dot at its upper right last.

3 + 2 + 4 = nine.

\subsection*{Your turn}
\begin{itemize}
\raggedright
  \item \emph{Write it:} \zh{突}, all nine, with the model in front of you
  \item \emph{Cover:} write it again from memory, and count the pieces, not the strokes
\end{itemize}

\begin{tcolorbox}[breakable,colback=teal!4,colframe=teal!35!black,title={Before you move on}]
How many strokes in \zh{突}? (\textbf{Nine}.) Which character you write has the same roof? (\textbf{\zh{家}}.) Where have you met the piece at the foot? (\textbf{On the right of \zh{然}}.)

Source: \href{https://raw.githubusercontent.com/chanind/hanzi-writer-data/68d10a4b21150cae5e1ebbd223eed289cf32d90c/data/\%E7\%AA\%81.json}{Hanzi Writer Data \zh{突}, pinned stroke paths and medians}.
\end{tcolorbox}
\section[\texorpdfstring{\zh{一下子} (yíxiàzi)}{一下子 (yíxiàzi)}]{\zh{一下子} (yíxiàzi) — all at once}

\label{lesson:ZH-C279-yixiazi}



\begin{quote}
Before the new one: last lesson's word, then \zh{一} with its rising tone.

\begin{itemize}
\raggedright
  \item \emph{Read it:} \zh{突然}, aloud — \emph{tūrán}
  \item \emph{Recall:} say \zh{一个} — \emph{yí ge}, with \zh{一} rising — \textbf{R3}, durable
  \item \emph{Recall:} say \zh{冷} — \emph{lěng}, cold — \textbf{R3}
\end{itemize}
\end{quote}

\subsection*{You'll want to know: \zh{一下子}}
\textbf{\zh{一下子}} — \emph{yíxiàzi} — \textquotedblleft{}all at once; in no time\textquotedblright{}.

\begin{quote}
\textbf{\zh{天气一下子冷了}} \emph{tiānqì yíxiàzi lěng le} — All at once the weather turned cold.
\end{quote}

\zh{一} rises before the fall of xià: \emph{yíxiàzi}, as in \emph{yí ge}.

Every character in it is one you write.

\begin{cousinweb}
\textbf{\zh{一下子}} (\emph{yíxiàzi}) is built from \zh{一} 'one', \zh{下}, here a single go, and the ending \zh{子}. All three are inherited from Middle Chinese, \emph{'jit}, \emph{haeX} and \emph{tsiX} (Guangyun, in Baxter's transcription). The dictionary lists \zh{下} under two readings, \emph{haeX} and \emph{haeH}, and glosses the first 'below' and 'to go down'; neither 'a go' nor 'all at once' is given there, so both are the modern word's own senses.
\end{cousinweb}

\begin{grammarlens}[title={What you've built}]
\zh{突然} says it came without warning; \zh{一下子} says the change took no time at all.

Both go before the verb: \textbf{\zh{天气一下子冷了}} \emph{tiānqì yíxiàzi lěng le}. The \zh{了} at the end is the \zh{了} of a change.
\end{grammarlens}

\subsection*{Your turn}
\begin{itemize}
\raggedright
  \item \emph{Say it:} \emph{yíxiàzi}
  \item \emph{Say it:} \emph{yíxiàzi}, once more
  \item \emph{Say it:} \zh{天气一下子冷了} — \emph{tiānqì yíxiàzi lěng le} — All at once the weather turned cold.
  \item \emph{Say it:} \zh{我一下子明白了} — \emph{wǒ yíxiàzi míngbai le} — All at once I understood.
  \item \emph{Recall:} say \zh{突然}, then \zh{一下子}
\end{itemize}

\begin{tcolorbox}[breakable,colback=teal!4,colframe=teal!35!black,title={Before you move on}]
What does \zh{一下子} mean? (\textquotedblleft{}All at once\textquotedblright{}.) How does \zh{一} sound in it? (\textbf{Rising: \emph{yí}}.)
\end{tcolorbox}
\section[(reading)]{(reading) — a day on the mountain}

\label{lesson:ZH-C279-reading}



\begin{quote}
Thirteen short lines: a day that went wrong, then right.

\begin{itemize}
\raggedright
  \item \emph{Recall:} say \zh{一下子} — \emph{yíxiàzi}, all at once — \textbf{R1}, one lesson back
\end{itemize}

\begin{quote}
\zh{山上} \zh{山下}
\end{quote}

\emph{Read it:} both aloud — \emph{shān shàng}, \emph{shān xià}
\end{quote}

\begin{tcolorbox}[breakable,skin=enhanced,colback=orange!6,colframe=orange!45!black,boxrule=0.5pt,arc=1mm,left=6pt,right=6pt,top=4pt,bottom=4pt,fonttitle=\bfseries,title={Read this}]
\begin{quote}
\zh{上个星期六} \zh{我和两个同学去山上} \zh{我们先开车} \zh{然后下车看花} \zh{中午我们在山上吃饭} \zh{吃了饭以后} \zh{突然下雨了} \zh{雨很大} \zh{我们都没有雨衣} \zh{我们正在找车的时候} \zh{一个老人来了} \zh{他说} \zh{去我家} \zh{他家在山下} \zh{我们在他家喝了茶} \zh{第二天} \zh{天气很好} \zh{我们先去看了山上的花} \zh{然后回家了} \zh{最后} \zh{我们给他打了电话} \zh{说谢谢} \zh{那天是我最开心的一天}
\end{quote}

\emph{Read it:} the story through once without stopping; do not translate as you go
\end{tcolorbox}

\begin{tcolorbox}[breakable,skin=enhanced,colback=orange!6,colframe=orange!45!black,boxrule=0.5pt,arc=1mm,left=6pt,right=6pt,top=4pt,bottom=4pt,fonttitle=\bfseries,title={Read this}]
\begin{itemize}
\raggedright
  \item \emph{shàng ge xīngqīliù, wǒ hé liǎng ge tóngxué qù shān shàng} — Last Saturday two classmates and I went up the mountain.
  \item \emph{wǒmen xiān kāichē, ránhòu xiàchē kàn huā} — First we drove, then we got out to look at the flowers.
  \item \emph{zhōngwǔ wǒmen zài shān shàng chī fàn} — At noon we ate on the mountain.
  \item \emph{chī le fàn yǐhòu, tūrán xiàyǔ le} — After we had eaten, it suddenly started to rain.
  \item \emph{yǔ hěn dà, wǒmen dōu méiyǒu yǔyī} — The rain was heavy, and none of us had a raincoat.
  \item \emph{wǒmen zhèngzài zhǎo chē de shíhou, yí ge lǎorén lái le} — While we were looking for the car, an old man came along.
  \item \emph{tā shuō, qù wǒ jiā} — He said: come to my house.
  \item \emph{tā jiā zài shān xià} — His house was at the foot of the mountain.
  \item \emph{wǒmen zài tā jiā hē le chá} — We had tea at his house.
  \item \emph{dì'èr tiān, tiānqì hěn hǎo} — The next day the weather was fine.
  \item \emph{wǒmen xiān qù kàn le shān shàng de huā, ránhòu huí jiā le} — First we went to see the flowers on the mountain, then we went home.
  \item \emph{zuìhòu, wǒmen gěi tā dǎ le diànhuà, shuō xièxie} — Last of all, we phoned him to say thank you.
  \item \emph{nà tiān shì wǒ zuì kāixīn de yì tiān} — That was my happiest day.
\end{itemize}
\end{tcolorbox}

\begin{tcolorbox}[breakable,skin=enhanced,colback=orange!6,colframe=orange!45!black,boxrule=0.5pt,arc=1mm,left=6pt,right=6pt,top=4pt,bottom=4pt,fonttitle=\bfseries,title={Read this}]
\emph{Read it:} the story a second time, the pinyin covered
\end{tcolorbox}

\subsection*{Your turn}
\begin{itemize}
\raggedright
  \item Put these in order: the tea, the rain, the drive, the phone call. (\textbf{The drive, the rain, the tea, the phone call}.)
  \item The next day? (\textbf{\zh{先去看了山上的花} \zh{然后回家了}}.)
  \item How did they thank him? (\textbf{\zh{给他打了电话}}.)
\end{itemize}

\begin{tcolorbox}[breakable,colback=teal!4,colframe=teal!35!black,title={Before you move on}]
Who came while they looked for the car? (\textbf{\zh{一个老人}}.)
\end{tcolorbox}
\section[Practice]{\zh{他突然来了}}

\label{lesson:ZH-C279-practice}



\begin{quote}
The chapter closes here: the story you read, told aloud as your own.

\begin{itemize}
\raggedright
  \item \emph{Read it:} \zh{第二天} \zh{天气很好}, aloud — \emph{dì'èr tiān, tiānqì hěn hǎo}
  \item \emph{Recall:} say \zh{一下子} — \emph{yíxiàzi} — \textbf{R2}
  \item \emph{Recall:} say \zh{我回家的时候} \zh{妈妈正在做饭} — \emph{wǒ huí jiā de shíhou, māma zhèngzài zuòfàn} — \textbf{R2}
\end{itemize}
\end{quote}

\subsection*{The exchange}
\begin{quote}
\textbf{A:} \zh{上个星期六你做什么了}
\end{quote}

\begin{quote}
\emph{shàng ge xīngqīliù nǐ zuò shénme le} — What did you do last Saturday?
\end{quote}

\begin{quote}
\textbf{B:} \zh{我和两个同学去山上了}
\end{quote}

\begin{quote}
\emph{wǒ hé liǎng ge tóngxué qù shān shàng le} — Two classmates and I went up the mountain.
\end{quote}

\begin{quote}
\textbf{A:} \zh{天气好吗}
\end{quote}

\begin{quote}
\emph{tiānqì hǎo ma} — Was the weather good?
\end{quote}

\begin{quote}
\textbf{B:} \zh{吃了饭以后} \zh{突然下雨了}
\end{quote}

\begin{quote}
\emph{chī le fàn yǐhòu, tūrán xiàyǔ le} — After we'd eaten, it suddenly rained.
\end{quote}

\begin{quote}
\textbf{A:} \zh{你们有雨衣吗}
\end{quote}

\begin{quote}
\emph{nǐmen yǒu yǔyī ma} — Did you have raincoats?
\end{quote}

\begin{quote}
\textbf{B:} \zh{没有} \zh{我们找车的时候} \zh{一个老人来了}
\end{quote}

\begin{quote}
\emph{méiyǒu, wǒmen zhǎo chē de shíhou, yí ge lǎorén lái le} — No. While we were looking for the car, an old man came.
\end{quote}

\begin{quote}
\textbf{A:} \zh{第二天呢}
\end{quote}

\begin{quote}
\emph{dì'èr tiān ne} — And the next day?
\end{quote}

\begin{quote}
\textbf{B:} \zh{天气一下子好了} \zh{我们看了花} \zh{然后回家了}
\end{quote}

\begin{quote}
\emph{tiānqì yíxiàzi hǎo le, wǒmen kàn le huā, ránhòu huí jiā le} — The weather was fine in no time. We looked at the flowers, then went home.
\end{quote}

B tells the story from the inside: \zh{我们}, not they. The time words still lead each turn.

\begin{tcolorbox}[breakable,skin=enhanced,colback=orange!6,colframe=orange!45!black,boxrule=0.5pt,arc=1mm,left=6pt,right=6pt,top=4pt,bottom=4pt,fonttitle=\bfseries,title={Read this}]
\emph{Read it:} the exchange through once without stopping, then answer the question below
\end{tcolorbox}

\begin{grammarlens}[title={What you've built — telling a day}]
A story needs a moment, a scene, a turn and a next day:

\begin{itemize}
\raggedright
  \item \textbf{\zh{他来的时候} \zh{我正在看书}} \emph{tā lái de shíhou, wǒ zhèngzài kàn shū} — the moment and the scene.
  \item \textbf{\zh{突然下雨了}} \emph{tūrán xiàyǔ le} — the turn.
  \item \textbf{\zh{天气一下子冷了}} \emph{tiānqì yíxiàzi lěng le} — a change that took no time.
  \item \textbf{\zh{第二天}} \emph{dì'èr tiān} — the day after, inside the story.
\end{itemize}

\textbf{Listen.} \emph{Nà tiān nǐ zuò shénme le?}

\textbf{Speak.} Tell a day of your own in six lines: a time, what you were doing, what happened, the next day.

\emph{Write it:} \zh{第二天} and \zh{突然}
\end{grammarlens}

\subsection*{Your turn}
\begin{itemize}
\raggedright
  \item \emph{Say it:} \zh{那天我正在看书的时候} \zh{他突然来了} — \emph{nà tiān wǒ zhèngzài kàn shū de shíhou, tā tūrán lái le} — That day, while I was reading, he suddenly came.
  \item \emph{Say it:} \zh{第二天天气很好} — \emph{dì'èr tiān tiānqì hěn hǎo} — The next day the weather was fine.
  \item \emph{Say it:} the whole exchange, both parts
  \item \emph{Write it:} \zh{突然下雨了}
\end{itemize}

\begin{tcolorbox}[breakable,colback=teal!4,colframe=teal!35!black,title={Before you move on}]
Say suddenly it rained. (\textbf{\zh{突然下雨了}}.) Say when he came, I was eating. (\textbf{\zh{他来的时候} \zh{我正在吃饭}}.) Say the next day. (\textbf{\zh{第二天}}.)
\end{tcolorbox}
